
% \subsection{AES}
% \label{sec:appendix_aes}
% \begin{table*}[h!]
%     \centering
%     \begin{tabular}{|p{2cm}|p{5cm}|p{5cm}|}
%         \hline
%          & ASAP & ASAP++ \\
%          \hline\hline
%         prompt 1 & elaboration, organization, fluency, \textbf{audience awareness} & content, organization, word choice, sentence fluency, \textbf{conventions}\\
%         \hline
%         prompt 4 & addresses the question, using information from the text, conclusion or understanding about the text & content, prompt adherence, \textbf{language}, \textbf{narrativity}\\
%         \hline
%         prompt 6 & relevance, \textbf{specific information} & content, prompt adherence, \textbf{language}, \textbf{narrativity}\\
%         \hline
%     \end{tabular}
%     \caption{Aspects mentioned by ASAP and ASAP++ in the rubrics. Bolded aspects indicate aspects that are mentioned in ASAP or ASAP++ but not both. }
%     \label{tab:asap_mismatch}
% \end{table*}

\subsection{Prompt for AES}
\label{sec:aes_prompt_appendix}
The holistic and analytic rubrics within the prompts are taken from ASAP \citep{asap-aes} and ASAP++ \citep{mathias-bhattacharyya-2018-asap}. The contextual information present in the grading guidelines are formatted into the prompt, and additionally added with slight edits to the analytic prompts in the edited condition. The examples chosen from ASAP++ are essay IDs 449, 1264, 1616, 9125, 9430, 9497, 15558, 16520, and 16581.

\begin{tcolorbox}[title=Prompt 1 Holistic Context,breakable]
Instructions:

1) The following essay is a first draft written by an 8th grade student in forty-five minutes in reaction to a prompt designed to elicit persuasive writing. You will score these timed responses holistically, which means that you will determine a score based on the overall impression most often gained from a single reading of the response.

2) You will be given a rubric that outlines a six-point scale. Each score point on that scale is described by an overall statement which captures the essence of the response. The elements of the response (elaboration, organization, fluency and audience awareness) that are typical for that score point are described below the overall statement. Individual responses may be stronger in one feature and weaker in another. In other words, the list of features at each score point, while helpful, cannot perfectly describe every response in a score point category.

3) A committee of expert readers uses this rubric as a guide to select anchor papers for each score point. Anchor papers are examples of actual student work. The committee prepares an anchor set composed of several papers at each score point. They deliberately select papers to show an appropriate range of writing skill for each score point and to represent the variety of approaches students take when addressing the writing prompt. You rely heavily on these anchor sets to guide your scoring.

4) Errors in spelling, punctuation, grammar, and usage are not considered as part of the criteria for scoring. Also, papers receive a score based on the work the student did complete even if they seem to be unfinished. Because the writing sample is a timed response, it is generally assumed that these errors and omissions could have been corrected if the student had been given an opportunity to revise and edit. You are trained to read through these errors when you score student papers.

5) Score the essay on a scale from 1 to 6.
\end{tcolorbox}

\begin{tcolorbox}[title=Prompt 4 Holistic Context,breakable]
    Instructions:
    
1) The following essay is written by an 10th grade student in response to a prompt that is dependent on reading the story provided.

2) You will be given a rubric that outlines a four-point scale.

3) Training materials consist of a rubric and a scoring guide of ten responses.

4) Score the essay on a scale from 0 to 3.
\end{tcolorbox}

\begin{tcolorbox}[title=Prompt 6 Holistic Context,breakable]
    Instructions:
    
1) The following essay is written by an 10th grade student in response to a prompt that is dependent on reading the excerpt provided.

2) You have a four year baccalaureate degree as well as documented coursework in English. You are not a teacher, substitute teacher, support staff, tutor, administrator, etc., who is currently under contract or employed by or in schools, or under 18 years of age.

3) You will be given a rubric that outlines a five-point scale.

4) You will be given an anchor set which will consist of responses that are typical, rather than unusual or uncommon; solid, rather than controversial or borderline; and true, meaning that these have scores that cannot be changed by anyone other than pertinent personnel. Anchor sets will typically have 2 to 3 sample responses at each score point (the middle score points will have 3 sample responses, 1 representing the mid-high to high end of the score point range, 1 in the middle, and 1 at the mid-low to low end).

5) Score the essay on a scale from 0 to 4.
\end{tcolorbox}

\begin{tcolorbox}[title=Example of context given to an analytic prompt for separate (Prompt 1),breakable]
Instructions:

1) The following essay is a first draft written by an 8th grade student in forty-five minutes in reaction to a prompt designed to elicit persuasive writing.

2) You will be given a rubric that outlines a six-point scale for an attribute. 

3) Score the essay on a scale from 1 to 6 on the attribute.
\end{tcolorbox}

\begin{tcolorbox}[title=Example of context given to an analytic prompt for edited (Prompt 1),breakable]
Instructions:

1) The following essay is a first draft written by an 8th grade student in forty-five minutes in reaction to a prompt designed to elicit persuasive writing.

2) You will be given a rubric that outlines a six-point scale for an attribute. 

3) A committee of expert readers uses this rubric as a guide to select anchor papers for some score points. Anchor papers are examples of actual student work. The committee prepares an anchor set composed of several papers at various score points. They deliberately select papers to show an appropriate range of writing skill and to represent the variety of approaches students take when addressing the writing prompt. You rely heavily on these anchor sets to guide your scoring.

4) Also, papers receive a score based on the work the student did complete even if they seem to be unfinished. Because the writing sample is a timed response, it is generally assumed that these errors and omissions could have been corrected if the student had been given an opportunity to revise and edit. You are trained to read through these errors when you score student papers.

5) Score the essay on a scale from 1 to 6 on the attribute.
\end{tcolorbox}
